\documentclass[11pt]{article}
\usepackage[a4paper,margin=25mm]{geometry}
\usepackage[T1]{fontenc}
\usepackage{lmodern,amsmath,amssymb,amsthm,booktabs,array}
\usepackage{microtype}
\usepackage[colorlinks=true,linkcolor=blue,citecolor=blue,urlcolor=blue]{hyperref}
\usepackage{listings,float}
\lstset{basicstyle=\ttfamily\footnotesize,columns=fullflexible,keepspaces=true,breaklines=true,frame=single}
\newtheorem{theorem}{Theorem}
\newtheorem{lemma}[theorem]{Lemma}
\newcommand{\Lfun}{\mathcal L}
\newcommand{\bb}{\boldsymbol}
\newcommand{\Q}{\mathbb Q}
\newcommand{\N}{\mathbb Z_{\geq0}}
\title{A four-moment proof of the conjectured\protect\\recurrence for A172671}
\author{GPT-6 Astra}
\date{30 October 2026}
\hypersetup{pdfauthor={GPT-6 Astra},pdftitle={A four-moment proof of the conjectured recurrence for A172671}}
\begin{document}
\maketitle
\begin{abstract}
Let $a_n$ count the $3n\times6$ matrices with entries in $\{0,1,2\}$,
row sums $2$, and column sums $n$, with $a_0=1$.
We prove the exact fourth-order recurrence conjectured by Kauers and
Koutschan for $b_n=(n!)^3a_n/(3n)!$, for every integer $n\geq0$.
A formal Gaussian moment functional converts the coefficient count into
the sixth moment of a modified Hermite polynomial. Integration by parts
closes on four moments and gives an explicit $4\times4$ transfer matrix.
A finite polynomial identity establishes the conjectured recurrence. Two exact
arithmetic implementations verify the identity, and the accompanying
code reproduces the certificate.
\end{abstract}

\section{Counting problem, prior work, and exact result}
Rows and columns are labeled. A row of sum $2$ is either $2$ in one
position, or $1$ in two distinct positions. There are $6+15=21$ choices.
Consequently, with
\[
 h(\bb t)=\sum_{1\leq i\leq j\leq6}t_it_j
 =\frac12\left(\left(\sum_{i=1}^6t_i\right)^2+\sum_{i=1}^6t_i^2\right),
\]
the exact counting formula is
\begin{equation}\label{eq:count}
 a_n=[t_1^n\cdots t_6^n]h(\bb t)^{3n},\qquad
 b_n=\frac{(n!)^3}{(3n)!}a_n.
\end{equation}
The convention $a_0=1$ counts the empty matrix. The OEIS entry
A172671 \cite{OEIS} starts at $n=1$; we explicitly include $n=0$.

Kauers and Koutschan \cite[Section 4.2, Conjecture 7]{KK} give the
recurrence below, and their Theorem 8 proves D-finiteness without
establishing this particular recurrence. Their final JIS article labels it a conjecture on page 20.
Yurkevich \cite[Conjecture 10.60 and equation (10.5), pp.~180--181]{Y}
also states the conjecture and uses the displayed decomposition of $h$
to compute terms more efficiently. That quadratic decomposition and the
recurrence itself are therefore prior work. The additional step here is
the explicit four-moment closure and its exact elimination certificate.

\newpage
\begin{theorem}\label{thm:main}
For every integer $n\geq0$, the sequence defined by \eqref{eq:count}
satisfies
\begin{equation}\label{eq:target}
 \sum_{j=0}^4 p_j(n)b_{n+j}=0,
\end{equation}
where the coefficients are exactly
\begin{align*}
 p_4(n)={}&3(n+3)(n+4)^3
 (3784n^4+32164n^3+100749n^2+137862n+69678),\\
 p_3(n)={}&-(n+3)\bigl(3799136n^7+72183584n^6+579689880n^5\\
 &\quad+2548427912n^4+6617561702n^3+10141503096n^2\\
 &\quad+8487349821n+2991586122\bigr),\\
 p_2(n)={}&-3\bigl(10844944n^8+222321352n^7+1973930222n^6\\
 &\quad+9916013134n^5+30831383530n^4+60768378830n^3\\
 &\quad+74160044251n^2+51243135187n+15352797306\bigr),\\
 p_1(n)={}&(n+2)\bigl(29681696n^7+504588832n^6+3602458816n^5\\
 &\quad+14001842392n^4+32010306742n^3+43078657918n^2\\
 &\quad+31639900193n+9799573455\bigr),\\
 p_0(n)={}&15435(n+1)^3(n+2)
 \bigl(3784n^4+47300n^3+219945n^2\\
 &\hspace{42mm}+450988n+344237\bigr).
\end{align*}
The initial conditions are
\begin{equation}\label{eq:initial}
 b_0=1,\qquad b_1=15,\qquad b_2=2249,\qquad b_3=444975.
\end{equation}
Since $p_4(n)>0$ for $n\geq0$, these four values and
\eqref{eq:target} determine the entire sequence.
\end{theorem}

\section{An exact Gaussian-moment representation}
Define a linear functional $\Lfun:\Q[x]\to\Q$ by
\begin{equation}\label{eq:functional}
 \Lfun(x^{2m})=(2m-1)!!=\frac{(2m)!}{2^m m!},\qquad
 \Lfun(x^{2m+1})=0,\qquad (-1)!!=1.
\end{equation}
It is the standard normal expectation on polynomials, but the proof
uses only the following algebraic identities:
\begin{equation}\label{eq:ibp}
 \Lfun(f')=\Lfun(xf),\qquad
 \Lfun(e^{tx})=e^{t^2/2}.
\end{equation}
For the first identity it suffices to take $f=x^r$: both sides vanish
when $r$ is even, while for $r=2m+1$ they are both $(2m+1)!!$.
The second identity follows coefficient by coefficient from
\eqref{eq:functional}. Exponential series below are formal series, and
$\Lfun$ is applied to their polynomial coefficients. No convergence or
interchange of infinite analytic integrals is required.

Define the modified Hermite polynomials by
\begin{equation}\label{eq:polynomials}
 e^{tx+t^2/2}=\sum_{n\geq0}P_n(x)\frac{t^n}{n!},\qquad
 P_n(x)=n!\sum_{k=0}^{\lfloor n/2\rfloor}
 \frac{x^{n-2k}}{2^k k!(n-2k)!}.
\end{equation}
The first polynomials are $1,x,x^2+1,x^3+3x$.
Differentiation of the generating series gives, for all $n\geq0$,
\begin{equation}\label{eq:hermite}
 P_{n+1}=xP_n+P_n',\qquad
 P_n''=nP_n-xP_n',\qquad P_{n+1}'=(n+1)P_n.
\end{equation}
For example, $\partial_x e^{tx+t^2/2}=t e^{tx+t^2/2}$ gives
$P_n'=nP_{n-1}$ for $n\geq1$, while
$\partial_t e^{tx+t^2/2}=(x+t)e^{tx+t^2/2}$ gives the first
relation; differentiating it gives the remaining identities.
At $n=0$ they hold directly since $P_0=1$.

\begin{lemma}\label{lem:moment}
For every $n\geq0$, $b_n=\Lfun(P_n^6)/(n!)^3$.
\end{lemma}
\begin{proof}
By \eqref{eq:ibp}, in the ring of formal series in $t_1,\ldots,t_6$,
\[
 \Lfun\left(\prod_{i=1}^6 e^{xt_i+t_i^2/2}\right)
 =\exp\left(\frac12\Bigl(\sum_i t_i\Bigr)^2+
             \frac12\sum_i t_i^2\right)=e^{h(\bb t)}.
\]
The coefficient of $t_1^n\cdots t_6^n$ on the left is
$\Lfun(P_n^6)/(n!)^6$. Since $h$ is homogeneous of degree $2$,
only $h^{3n}/(3n)!$ contributes to that coefficient on the right,
which is $a_n/(3n)!$. Multiplying by $(n!)^3$ proves the claim.
\end{proof}

\section{Closure on four moments}
Fix $n\geq0$, and put $y=P_n(x)$ and $z=P_n'(x)$.
Then $y'=z$ and $z'=ny-xz$. For $0\leq j\leq6$ write
\[
 e_j=y^{6-j}z^j,\qquad M_{k,j}=\Lfun(x^ke_j)\quad(k\geq0).
\]
Differentiating $x^ke_j$ and using \eqref{eq:ibp} gives
\begin{equation}\label{eq:reduction}
 (j+1)M_{k+1,j}=kM_{k-1,j}+(6-j)M_{k,j+1}+jnM_{k,j-1}.
\end{equation}
The first term is omitted at $k=0$. At $j=0$ and $j=6$, respectively,
the out-of-range term is absent because its coefficient is zero.
Explicitly,
\[
 (x^ke_j)'=kx^{k-1}e_j+(6-j)x^ke_{j+1}
                 +jnx^ke_{j-1}-jx^{k+1}e_j,
\]
so the factor $j+1$ in \eqref{eq:reduction} includes the contribution
from the Gaussian weight. In an integral formulation the boundary term
is a polynomial times $e^{-x^2/2}$, hence vanishes at both infinities;
the algebraic proof already accounts for it exactly.

By \eqref{eq:polynomials}, $P_n(-x)=(-1)^nP_n(x)$, so $e_j$ has
parity $j$. In particular $M_{0,j}=0$ for odd $j$.
Induction on $k$ in \eqref{eq:reduction} expresses every $M_{k,j}$
as a $\Q[n]$-linear combination of the four entries of
\[
 W_n=\bigl(\Lfun(y^6),\Lfun(y^4z^2),\Lfun(y^2z^4),
                  \Lfun(z^6)\bigr)^{\mathsf T}.
\]
All divisions in this induction are by the positive integers $j+1$;
there are no exceptional values of $n$.

\begin{lemma}\label{lem:transfer}
For every integer $n\geq0$,
\begin{equation}\label{eq:transfer}
 W_{n+1}=A(n)W_n,\qquad
 U_{n+1}=B(n)U_n,\quad
 U_n=\frac{W_n}{(n!)^3},\quad B(n)=\frac{A(n)}{(n+1)^3},
\end{equation}
where
\begin{equation}\label{eq:matrix}
\setlength{\arraycolsep}{4pt}
 A(n)=\begin{pmatrix}
 3(2n+1)(29n^2+29n+5)&15(76n^2+95n+28)&75(4n+3)&7\\[2pt]
 \dfrac{(n+1)^2(76n^2+57n+9)}3&\dfrac{241(n+1)^2(2n+1)}3&35(n+1)^2&0\\[6pt]
 (n+1)^4(4n+1)&21(n+1)^4&0&0\\[2pt]
 (n+1)^6&0&0&0
 \end{pmatrix}.
\end{equation}
The first component of $U_n$ is $b_n$, and $U_0=(1,0,0,0)^{\mathsf T}$.
\end{lemma}
\begin{proof}
Under $n\mapsto n+1$, \eqref{eq:hermite} replaces
$(y,z)$ by $(xy+z,(n+1)y)$. Thus, for $j=0,2,4,6$,
\begin{equation}\label{eq:entryrecipe}
 \Lfun(P_{n+1}^{6-j}(P_{n+1}')^j)
 =(n+1)^j\sum_{\ell=0}^{6-j}\binom{6-j}{\ell}
 M_{6-j-\ell,\ell}.
\end{equation}
Applying \eqref{eq:reduction} produces the corresponding row of
\eqref{eq:matrix}. For completeness, an exact finite specification of
this computation is as follows. Let $R_{k,j}$ be a row of length four,
set $R_{0,2r}$ equal to the $r$th unit row ($r=0,1,2,3$), and set
$R_{0,j}=0$ for odd $j$. Set out-of-range rows to zero. For $k\geq1$,
\begin{equation}\label{eq:R}
 R_{k,j}=\frac{(k-1)R_{k-2,j}+(6-j)R_{k-1,j+1}
                         +jnR_{k-1,j-1}}{j+1},
\end{equation}
omitting the first term when $k=1$. Row $r$ of $A(n)$ is then
\begin{equation}\label{eq:Arecipe}
 (n+1)^{2r}\sum_{\ell=0}^{6-2r}\binom{6-2r}{\ell}
 R_{6-2r-\ell,\ell}.
\end{equation}
Equations \eqref{eq:R}--\eqref{eq:Arecipe} verify all sixteen displayed
entries using only $k\leq6$. For instance,
$M_{2,0}=M_{0,0}+6M_{1,1}$ and
$M_{1,1}=(5M_{0,2}+nM_{0,0})/2$. Hence the row for $j=4$ is
$(n+1)^4\bigl((4n+1)M_{0,0}+21M_{0,2}\bigr)$, as displayed.
Normalization gives $B(n)$. The final assertions follow from
Lemma~\ref{lem:moment} and $P_0=1$, $P_0'=0$.
\end{proof}

\section{An exact elimination certificate}
Let $\epsilon=(1,0,0,0)$, and define matrix polynomials
\[
 C_0(n)=I_4,\qquad C_j(n)=A(n+j-1)C_{j-1}(n)\quad(1\leq j\leq4).
\]
The order of multiplication matters: explicitly,
$C_j(n)=A(n+j-1)\cdots A(n)$.

\begin{lemma}\label{lem:certificate}
With the five polynomials in Theorem~\ref{thm:main}, the following is
an identity of polynomial rows over $\Q[n]$:
\begin{equation}\label{eq:certificate}
 \boxed{\quad
 \sum_{j=0}^4 p_j(n)\left(\prod_{t=j+1}^{4}(n+t)^3\right)
                      \epsilon C_j(n)=(0,0,0,0).
 \quad}
\end{equation}
An empty product is $1$.
\end{lemma}
\begin{proof}
This is a finite polynomial identity with all of its inputs explicitly
given in \eqref{eq:matrix} and Theorem~\ref{thm:main}. It can be verified
by polynomial multiplication and coefficient collection as follows.
Starting with $C=I_4$ and $H=(0,0,0,0)$, for $j=0,\ldots,4$ replace
$C$ by $A(n+j-1)C$ if $j>0$, and add
$p_j(n)\prod_{t=j+1}^4(n+t)^3$ times the first row of $C$ to $H$.
The resulting polynomial coefficients are all zero. The sizes of the
four coordinate checks are given in Table~\ref{tab:certificate}; these
are identities in the indeterminate $n$, not evaluations at sample indices.

\begin{table}[H]
\centering
\begin{tabular}{lrrrr}
\toprule
Coordinate & 0 & 1 & 2 & 3\\
\midrule
Maximum degree of any summand &20&19&18&17\\
Coefficients checked (including constant)&21&20&19&18\\
Nonzero coefficients after summation&0&0&0&0\\
\bottomrule
\end{tabular}
\caption{Exact rational-polynomial certificate: all 78 coefficients vanish.}
\label{tab:certificate}
\end{table}

The calculation has been checked with two separate exact-arithmetic
programs. One uses SymPy rational functions; the other implements
polynomial addition, multiplication, substitution $n\mapsto n+h$, and
coefficient collection from scratch using rational numbers. The latter
also reconstructs \eqref{eq:matrix} from \eqref{eq:R}--\eqref{eq:Arecipe}
and reads the five target polynomials directly from the original TeX
excerpt. Thus neither recurrence guessing nor numerical tolerance enters
the certificate. Appendix~\ref{app:algorithm} gives precise coefficient
rules sufficient to reproduce the calculation using rational arithmetic alone.
\end{proof}

\begin{proof}[Proof of Theorem~\ref{thm:main}]
Iterating \eqref{eq:transfer} gives
\[
 b_{n+j}=\frac{\epsilon C_j(n)U_n}{\prod_{t=1}^j(n+t)^3}
 \qquad(0\leq j\leq4).
\]
For $n\geq0$, the common denominator
$D(n)=\prod_{t=1}^4(n+t)^3$ is strictly positive. Multiply the desired
recurrence by $D(n)$, substitute the last display, and apply
\eqref{eq:certificate} to $U_n$. This proves the exact recurrence on the
full claimed domain, including $n=0$.

The initial values can be computed without using that recurrence:
\[
 \Lfun(1)=1,\quad \Lfun(x^6)=15,\quad
 \Lfun((x^2+1)^6)=17992,\quad
 \Lfun((x^3+3x)^6)=96114600.
\]
Division by $(n!)^3$, for $n=0,1,2,3$, gives
\eqref{eq:initial}. Every factor of $p_4(n)$ outside its quartic is
positive, and all coefficients of that quartic are positive. Thus no
singular starting index occurs for $n\geq0$, and forward induction proves
the uniqueness assertion.
\end{proof}

\section{Independent checks and reproducibility}
The proof above uses exact algebra for unbounded $n$. Separately, a direct
coefficient dynamic program checks the original matrix interpretation.
If $F_r(\bb c)$ counts ordered lists of $r$ admissible rows with column
sums $\bb c$, then
\[
 F_0(\bb0)=1,\qquad
 F_{r+1}(\bb c)=\sum_{\bb s\in S}F_r(\bb c-\bb s),\qquad
 a_n=F_{3n}(n,n,n,n,n,n),
\]
where $S$ consists of the 21 row choices and any negative coordinate
gives zero. This independently agrees with \eqref{eq:polynomials} and
Lemma~\ref{lem:moment} for $0\leq n\leq8$. Further exact sixth-moment
computations through $n=40$ give zero recurrence residual for
$0\leq n\leq36$. These finite tests corroborate the proof; they are not
its justification.

\begin{table}[ht]
\centering
\begin{tabular}{rrr}
\toprule
$n$&$a_n$&$b_n$\\
\midrule
0&1&1\\
1&90&15\\
2&202410&2249\\
3&747558000&444975\\
4&3536978063850&102077289\\
5&19292117692187340&25493181015\\
\bottomrule
\end{tabular}
\caption{Exact examples, with the empty-matrix convention at $n=0$.}
\end{table}

The accompanying package contains the source excerpt, symbolic derivation,
independent rational-polynomial verifier, direct-count checks, and their
JSON reports. From the reproducibility directory the following commands reproduce the
checks:
\begin{lstlisting}
python src/verify_certificate.py
python src/route_alt_transfer.py
python src/independent_checks.py
\end{lstlisting}
The first command needs only the Python standard library. The second
uses SymPy; the third uses the standard library. Paths in the scripts
are resolved relative to their own locations. The target coefficients
are read from the supplied mathematical excerpt, whose transcription
was checked against the original JIS source. The verifier also compares
them to the coefficients printed in this paper. No network access is required.
The paper can be rebuilt with two runs of
\begin{lstlisting}
pdflatex -interaction=nonstopmode -halt-on-error -output-directory=output paper/main.tex
\end{lstlisting}
Rename the resulting \texttt{output/main.pdf} to
\texttt{output/research\_paper.pdf}, or use the supplied
\texttt{src/build\_paper.ps1} helper. A clean-directory release check
copies the indispensable inputs into a temporary subdirectory, repeats
the three verifiers and compilation, and records its results.

\paragraph{Authorship and verification.}
GPT-6 Astra is an AI system. The proof and computations were checked
by internal review and the exact verifiers described above; they have
not been externally peer reviewed or formalized in a proof assistant.

\appendix
\section{Coefficient arithmetic for the certificate}\label{app:algorithm}
Here is a complete finite arithmetic specification of the identity check.
Represent a polynomial $f$ by its list of rational coefficients
$[f_0,\ldots,f_d]$, with omitted coefficients interpreted as zero.
Use
\begin{align*}
 [n^k](f+g)&=f_k+g_k,\\
 [n^k](fg)&=\sum_{i=0}^{k}f_i g_{k-i},\\
 [n^k]f(n+h)&=\sum_{i=k}^{d}\binom{i}{k}h^{i-k}f_i.
\end{align*}
Apply addition and multiplication entrywise in the usual matrix-product
rule. Expand the factored entries in \eqref{eq:matrix} and the five
polynomials of Theorem~\ref{thm:main} by these rules, retaining exact
rational numbers throughout. The following pseudocode fixes all loop
bounds and the multiplication order:
\begin{lstlisting}
C := identity_matrix(4)
H := zero_row(4)
for j = 0,1,2,3,4:
    if j > 0:
        C := shift(A, j-1) * C
    q := p[j]
    for t = j+1,...,4:
        q := q * (n+t)^3
    H := H + q * first_row(C)
assert H[0] == 0 and H[1] == 0 and H[2] == 0 and H[3] == 0
\end{lstlisting}
In the final line equality means equality of every rational coefficient.
No interpolation, inferred annihilator, floating-point arithmetic, or
unproved closure assertion is used. Together with the displayed
polynomials and matrix, this algorithm fully specifies the 78 exact
zero-coefficient checks of Lemma~\ref{lem:certificate}.

\clearpage
\begin{thebibliography}{9}
\raggedright
\bibitem{KK} M. Kauers and C. Koutschan,
\emph{Some D-finite and some possibly D-finite sequences in the OEIS},
Journal of Integer Sequences \textbf{26} (2023), Article 23.4.5,
Section 4.2, Conjecture 7 and Theorem 8.
\url{https://cs.uwaterloo.ca/journals/JIS/VOL26/Koutschan/kout4.pdf}.
Preprint: \url{https://arxiv.org/abs/2303.02793v2}, 24 April 2023.
\bibitem{OEIS} The OEIS Foundation Inc., \emph{Entry A172671},
\url{https://oeis.org/A172671}, accessed 30 October 2026.
\bibitem{Y} S. Yurkevich, \emph{Integer sequences, algebraic series and differential operators},
doctoral dissertation, University of Vienna, 2023,
Conjecture 10.60 and equation (10.5), pp.~180--181.
\url{https://homepage.univie.ac.at/herwig.hauser/Publications/yurkevich%20dissertation%202023.pdf}.
\end{thebibliography}
\end{document}
